\section{Introduction}

Formal epistemology often begins with a state: a set of beliefs, a probability function, a Kripke model, or a theory closed under some consequence relation. The present project begins one step later. Its central question is not merely which epistemic states are acceptable, but which \emph{dynamics of change} make epistemic recovery possible when information is inconsistent, probabilistic, partially observed, and generated by mechanisms whose reliability is itself uncertain.

4PEL provides a useful environment for this question because contradiction is represented without immediate trivialization. An agent can occupy states in which positive and negative support coexist, while modal and probabilistic structure can still evolve. This separates phenomena that classical consistency-first accounts often place under one heading: being inconsistent now, being able to recover later, selecting which information to acquire, deciding whether to trust a sensor, and recognizing that the current model class may itself be inadequate.

The main claim of this paper is deliberately dynamic:

\begin{quote}
\textbf{Epistemic rationality can be studied as a property of learning, experimentation, calibration, and model-revision dynamics rather than solely as a property of instantaneous belief states.}
\end{quote}

The claim is not introduced as a philosophical stipulation. It is motivated by a sequence of machine-checked finite results. The development begins with Recovery variation, asymptotic stabilization, path dependence, and hysteresis. It isolates a confluent fragment of static evidence and a genuinely hysteretic regime generated by endogenous evidence. A state-sufficiency theorem then shows that history can matter through the current full model without requiring an additional history variable. Closed epistemic loops can return to the same Recovery status while leaving a nonzero full-model displacement, and oppositely composed anchored loops can fail to commute.

The second stage turns updates into actions. Once information acquisition is under agent control, \emph{being recovered} and \emph{learning well} come apart. A myopic Recovery reward can favor information avoidance, whereas a longer horizon can make temporary destabilization rational. Finite Bellman recursion, stochastic transition kernels, value of information, regret, and reward thresholds are then layered over the 4PEL dynamics. Partial observation adds a posterior over hidden epistemic states, and noisy observation makes Recovery a probabilistic signal rather than an oracle.

The third stage, new in v0.3.0, makes the experimental design itself endogenous. A controller may choose among observation channels with different likelihood profiles and costs. This exposes a fundamental distinction between \emph{more data} and \emph{more identifiability}: if two hidden states are observationally aliased under an experiment, arbitrarily many repetitions of that same experiment can leave the posterior split exactly unchanged. An alternative intervention can break the alias. The controller must therefore decide not only how long to observe, but \emph{what kind of observation to perform} and when additional information is no longer worth its cost.

The fourth and fifth stages turn this machinery reflexive. Posterior support is mapped back into the four values $T,F,B,N$, so contradiction and ignorance become control-relevant states rather than mere labels. Then uncertainty is lifted from the world to the observation model itself. An agent can be uncertain whether its sensor is reliable, update that uncertainty through calibration, and allow the resulting model posterior to change its next action. A joint posterior over world state and sensor model couples first-order and second-order uncertainty in one finite Bayesian state.

The final stage crosses the closed-model boundary. If every known model predicts an observation poorly, posterior reweighting within the known class may be the wrong operation. The formal development therefore introduces predictive surprise, model doubt, an explicit action to expand the hypothesis space, and finally an ``unknown'' probability mass. This yields a nested two-axis epistemic state
\[
  (\text{world status},\text{model status}) \in \{T,F,B,N\}^2,
\]
and a finite self-trust controller whose priority depends on the model-status coordinate. A trusted model licenses ordinary world-directed action; an ignorant model status triggers model expansion; contradictory or rejected model evidence triggers calibration.

Accordingly, the paper now studies a hierarchy of state descriptions:
\[
M_t
\quad\leadsto\quad
s_t
\quad\leadsto\quad
b_t
\quad\leadsto\quad
b_t^{\mathrm{world}\times\mathrm{model}}
\quad\leadsto\quad
(v^{\mathrm{world}}_t,v^{\mathrm{model}}_t).
\]
The first level carries 4PEL semantics, the second supports controlled hidden state, the third supports Bayesian control under partial observability, the fourth represents coupled uncertainty about world and sensor model, and the fifth compresses first- and second-order epistemic assessment into two four-valued coordinates. These levels are deliberately not identified with one another.

The principal contribution is therefore not a single universal optimality theorem. It is an executable and verified architecture that makes several boundaries explicit: Recovery versus truth, posterior uncertainty versus contradiction, repeated evidence versus identifiability, Bayesian optimism versus robust control, calibration versus world-directed sensing, and ordinary Bayesian updating versus structural model revision. The finite witnesses are strong enough to establish the possibility and internal consistency of these distinctions while leaving their generalization as an open mathematical program.

\section{Background and Scope}

\subsection{Four-valued epistemic states}

4PEL inherits the central insight of Belnap--Dunn four-valued semantics: positive and negative support need not be complements. A proposition may be supported positively, supported negatively, supported in both directions, or supported in neither direction. This supplies a paraconsistent semantic substrate in which inconsistency can be represented locally without validating arbitrary explosion \cite{belnap1977useful,dunn1976}.

The present paper does not reconstruct the complete 4PEL architecture. It uses only the following structural facts from the wider formalization:

\begin{itemize}
  \item formulas are evaluated in four-valued epistemic models;
  \item models contain probabilistic structure used to evaluate belief;
  \item evidence formulas can be used to Conditionalize the current model when admissibility conditions are satisfied;
  \item a compositional \emph{Recovery} predicate records when the relevant four-valued semantics supports a classical-like fragment for a target formula or modal expression.
\end{itemize}

Recovery is therefore a formal semantic property internal to the model. It must not be identified with correspondence truth. A recovered agent may still be factually wrong; conversely, a non-recovered state may encode genuine conflicting information about the world. The object studied here is the \emph{dynamics of Recovery}, not a general theory of truth.

\subsection{Updates and endogenous evidence}

Let $M$ be a current model and let $P$ be an admissible evidence formula. We write
\[
\Cond{P}(M)
\]
for the result of Conditionalization on the positive extension of $P$ in $M$.

A crucial distinction concerns how the relevant evidence event is obtained. For an atomic or otherwise belief-free formula, its extension is fixed by valuation and Boolean structure and is unaffected by a probability-only Conditionalization. By contrast, an evidence formula containing an epistemic belief operator can be \emph{endogenous}: its extension may depend on the probability distribution of the current model. Thus an earlier update can change what a later piece of syntactically identical evidence denotes.

This distinction separates the static and feedback regimes studied below.

\subsection{Formal method}

The results reported here were developed theorem-first in Lean 4. Countermodels and positive theorems are represented in the same environment. The formal workflow is deliberately adversarial: each proposed generalization is tested against explicit finite witnesses, and each philosophical interpretation is restricted to what the formal statement establishes.

The current snapshot is the 4PEL-Lean-Formalization repository at commit\
\begin{center}
\small\path{67f1187d4b64f77c43289999c4e22b004a3dafa3}
\end{center}
The principal Lean modules for this paper are listed in Appendix A; they include the path-dependence, hysteresis, semantic-feedback, confluence, state-sufficiency, and tail-attractor developments \cite{repo2026}.

\section{From Recovery to Path-Dependent Dynamics}

\subsection{Finite variation is not attraction}

The early Recovery-dynamics results introduce a natural-valued potential $\Phi(s)$ on epistemic states. Under a basic descent contract, a change in Boolean Recovery status must consume potential. The immediate consequence is a finite variation principle: Recovery cannot flip infinitely often along a valid trajectory whose relevant potential is finite.

This already separates two questions:
\[
\text{Does the phase eventually stop changing?}
\qquad\text{and}\qquad
\text{Which phase is eventually reached?}
\]
A finite flip budget answers only the first. A trajectory may stabilize permanently in Recovery or permanently in Non-Recovery.

This distinction is methodologically important. Convergence of an observable is not yet convergence to an epistemically privileged observable value.

\subsection{Path dependence from one recovered start}

\Gate{29} establishes a stronger result. There exists a single recovered 4PEL model with two admissible evidence futures whose asymptotic Recovery phases differ. In the concrete witness, learning one atomic proposition selects a fixed point that remains recovered; learning another selects a fixed point that remains non-recovered. Repeating the same admissible atomic evidence is idempotent after the first update.

\begin{proposition}[Path-dependent Recovery phase]
There is a recovered model $M_0$ and two admissible update schedules $\sigma_1,\sigma_2$ such that the corresponding trajectories begin at the same $M_0$, while one is eventually permanently recovered and the other eventually permanently non-recovered.
\end{proposition}

The proposition separates three notions:
\begin{enumerate}
  \item Recovery now;
  \item eventual Recovery along one admissible future;
  \item Recovery under all admissible futures.
\end{enumerate}
The first does not imply the third. Recovery is therefore not, by itself, a guarantee about the future selected by later evidence.

\subsection{Recovery as an epistemic attractor}

\Gate{32} asks what additional structure is sufficient to select Recovery as the eventual phase. The answer uses two local conditions on a descent system:

\begin{align*}
\text{Recovery absorption:}\quad
&\Recovery(s)\land s\to t \Longrightarrow \Recovery(t),\\[2mm]
\text{Non-Recovery consumption:}\quad
&\neg\Recovery(s)\land s\to t \Longrightarrow \Phi(t)<\Phi(s).
\end{align*}

\begin{theorem}[Finite Recovery attraction]
If Recovery is absorbing and every Non-Recovery transition strictly decreases a natural-valued potential, then every valid trajectory reaches permanent Recovery in at most $\Phi(s_0)$ steps:
\[
\exists N\leq \Phi(s_0)\;\forall n\geq N:\;\Recovery(s_n).
\]
\end{theorem}

The potential now has a stronger interpretation than a mere flip budget. Under the attractor contract it behaves as a distance-like resource: remaining outside Recovery costs finite epistemic capacity.

The theorem is conditional. Arbitrary 4PEL Conditionalization does not satisfy Recovery absorption, as the path-dependence witness itself demonstrates. The result instead identifies a structural contract under which Recovery becomes a quantitative attractor.

\section{Hysteresis from Endogenous Evidence}

\subsection{The Gate-33 witness}

Path dependence does not yet amount to hysteresis. A stronger question holds the evidence ingredients fixed and changes only their order. \Gate{33} supplies such a witness with the pair $\{e,\Bel p\}$.

Initially the positive extension of $\Bel p$ contains all three relevant worlds. After conditioning on $e$, the belief profile changes, and the positive extension of the \emph{same syntactic evidence formula} $\Bel p$ shrinks. Consequently,
\[
\Cond{e}(\Cond{\Bel p}(M_0))
\quad\text{and}\quad
\Cond{\Bel p}(\Cond{e}(M_0))
\]
have different Recovery outcomes.

\begin{proposition}[Endogenous-evidence hysteresis]
There is a model $M_0$ and two admissible evidence formulas $e$ and $\Bel p$ such that reversing their order changes the final Recovery judgment:
\[
\neg\Recovery(\Bel p;e)
\qquad\text{but}\qquad
\Recovery(e;\Bel p).
\]
\end{proposition}

The result must be interpreted carefully. It does \emph{not} show that two fixed extensional events fail to commute. Rather, the evidence language can inspect the evolving epistemic model, and the first update changes the event denoted by the second formula. The source of hysteresis is therefore semantic feedback.

\subsection{Regenerative Recovery budgets}

Hysteresis motivates a second relaxation. If epistemic updates can reveal or create new relevant structure, a fixed initial flip budget may be too restrictive. \Gate{34} introduces a regeneration term $\rho(s,t)$ and replaces pure descent with a balance law:
\[
 c(s,t)+\Phi(t)\leq \Phi(s)+\rho(s,t),
\]
where $c(s,t)$ is the unit cost of a Recovery-status change.

Along a finite path this telescopes to
\[
\#\text{Recovery changes}+\Phi_{\mathrm{final}}
\leq
\Phi_{0}+\sum \rho.
\]
Hence
\[
\boxed{
\#\text{Recovery changes}
\leq
\Phi_{0}+\text{total regeneration}.
}
\]

The zero-regeneration case recovers the earlier finite-descent theory. The conceptual shift is from a depleting counter to a balance sheet. A Recovery transition that could not be financed by the original potential may still occur if earlier updates regenerate resource.

Unbounded regeneration destroys the stabilization guarantee. This observation becomes decisive in the open problem of growing ontologies discussed in Section~\ref{sec:open}.

\section{The Confluence Boundary}

\subsection{Belief-free evidence is extension-stable}

\Gate{36} defines a recursively belief-free fragment of the evidence language. Conditionalization changes the probability measure but leaves the atomic valuation and accessibility structure fixed. Structural induction therefore yields exact evaluation stability for belief-free formulas. Their positive evidence extensions do not change under prior Conditionalization.

In schematic form,
\[
\mathsf{BeliefFree}(P)
\Longrightarrow
\llbracket P\rrbracket^+_{\Cond{Q}(M)}
=
\llbracket P\rrbracket^+_M.
\]

This gives a syntactically identifiable region in which the Gate-33 feedback mechanism cannot arise.

\subsection{Static schedule confluence}

\Gate{38} lifts this stability from single formulas to finite schedules. For arbitrary belief-free formulas $P$ and $Q$, under finite probability integrity and the required sequential admissibility conditions,
\[
\Cond{P}(\Cond{Q}(M))
\ObsEq
\Cond{Q}(\Cond{P}(M)),
\]
where $\ObsEq$ denotes modal observational equivalence.

Moreover, observational equivalence is preserved under a common admissible belief-free suffix. Thus an adjacent swap inside a finite static schedule preserves every modal value and every derived compositional Recovery judgment, provided both reordered executions remain admissible.

This yields a local generator for permutation confluence. It also reveals the beginnings of an update algebra: atomic repetition is idempotent, and belief-free updates commute observationally under the stated conditions.

\begin{figure}[ht]
\centering
\begin{tikzpicture}[
  node distance=12mm,
  big/.style={draw=black!55, rounded corners=3pt, fill=boxgray, align=center, minimum width=58mm, minimum height=22mm, inner sep=7pt},
  arr/.style={-{Latex[length=2.2mm]}, semithick, draw=black!65}
]
\node[big] (static) {\textbf{Extension-stable evidence}\\belief-free is a sufficient fragment\\$P;Q\;\ObsEq\;Q;P$ under admissibility};
\node[big, right=18mm of static] (feedback) {\textbf{Extension-sensitive evidence}\\semantic feedback is possible\\Recovery hysteresis becomes possible};
\draw[arr] ($(static.east)+(0,3mm)$) -- node[above, font=\small]{boundary} ($(feedback.west)+(0,3mm)$);
\draw[arr] ($(feedback.west)+(0,-4mm)$) -- node[below, font=\small]{not an iff} ($(static.east)+(0,-4mm)$);
\end{tikzpicture}
\caption{The verified phase boundary. Belief-free syntax guarantees extension stability, but belief-containing syntax does not by itself imply sensitivity.}
\label{fig:boundary}
\end{figure}

\subsection{A semantic no-go theorem for hysteresis}

\Gate{40} removes the remaining syntactic ambiguity. Let $P$ and $Q$ be arbitrary evidence formulas, possibly containing belief operators. Suppose their positive extensions are mutually stable under the opposite update. Under probability integrity and mutual sequential admissibility, the two update orders are modally observationally equivalent and therefore Recovery-confluent.

\begin{theorem}[Strong hysteresis boundary]
If $P$ is extension-stable under the update by $Q$ and $Q$ is extension-stable under the update by $P$, then
\[
P;Q\ObsEq Q;P,
\]
and all current modal Recovery judgments agree at the two endpoints.
\end{theorem}

Taking the contrapositive yields the central diagnostic:
\[
\boxed{
\text{Recovery hysteresis}
\Longrightarrow
\text{semantic extension feedback in at least one direction}.
}
\]

Thus the presence of a belief operator is not the explanatory primitive. A belief-containing formula may remain extension-stable in a particular model. The verified necessary condition is semantic, not merely syntactic. Feedback is necessary for the formal two-update Recovery hysteresis result, but it is not sufficient by itself.

\section{History Sensitivity without a History Register}

Hysteresis naturally raises a memory question. If update order matters, must the formal state contain an explicit record of the past?

\Gate{39} answers negatively for Conditionalization. A finite schedule-run object allows arbitrary evidence formulas, including belief-dependent evidence. The next model is determined by the current full model, the next syntactic evidence formula, and an admissibility proof. Proof irrelevance prevents different admissibility witnesses from encoding hidden dynamics.

\begin{theorem}[Epistemic state sufficiency]
For arbitrary admissible Conditionalization schedules,
\[
\boxed{
\text{same current full model}
+
\text{same future evidence schedule}
\Longrightarrow
\text{same endpoint full model}.
}
\]
\end{theorem}

Two histories that merge at exactly the same current model therefore coalesce under every common admissible future schedule. Since endpoint model equality implies agreement on subsequent modal values and Recovery judgments, the full current model is a sufficient statistic for the represented update dynamics.

This reconciles path dependence with a Markov-style interpretation:
\[
\text{history}
\longrightarrow
\text{different present model}
\longrightarrow
\text{different future semantics}.
\]
No additional memory variable is required. The past matters precisely insofar as it has become present structure.

The qualification \emph{full model} is essential. Recovery status alone is too coarse. Two states can share the same current Boolean Recovery phase while encoding probability distributions or belief extensions that make their future behavior diverge.
